\documentclass[10pt,a4paper]{article}

% --- XeLaTeX: fonts + Russian ---
\usepackage{fontspec}
\usepackage{polyglossia}
\setdefaultlanguage{russian}
\setotherlanguage{english}
\setmainfont{PT Sans}
\setmonofont[Scale=MatchLowercase]{PT Mono}
\newfontfamily\cyrillicfonttt[Scale=MatchLowercase]{PT Mono}

% --- layout ---
\usepackage[a4paper,top=1.0cm,bottom=1.0cm,left=1.5cm,right=1.5cm]{geometry}
\usepackage{titlesec}
\usepackage{enumitem}
\usepackage{xcolor}
\usepackage[hidelinks]{hyperref}

\definecolor{rule}{HTML}{333333}

% section headers: bold, small rule under
\titleformat{\section}{\normalsize\bfseries}{}{0em}{\MakeUppercase}[{\color{rule}\titlerule[0.6pt]}]
\titlespacing*{\section}{0pt}{6pt}{3pt}

% tight bullet lists
\setlist[itemize]{leftmargin=1.15em, labelsep=0.5em, itemsep=1.6pt, topsep=2pt, parsep=0pt, before=\vspace{1pt}}
\renewcommand{\labelitemi}{\textbullet}

\setlength{\parindent}{0pt}
\setlength{\parskip}{0pt}
\pagestyle{empty}
\linespread{1.0}
\tolerance=1500
\emergencystretch=2.5em

% role line: bold title left, date right
\newcommand{\role}[2]{\textbf{#1}\hfill #2\par}

\begin{document}

% ===== Header =====
\begin{center}
  {\LARGE\bfseries Михайловский Илья Евгеньевич}\\[2pt]
  {\large Python Backend Developer}\\[3pt]
  Москва \textperiodcentered{} удалённо или гибрид\\[2pt]
  Telegram: \href{https://t.me/Iluuuaaa}{@Iluuuaaa} \textperiodcentered{}
  \href{mailto:ilyamihailovsy@gmail.com}{ilyamihailovsy@gmail.com} \textperiodcentered{}
  \href{https://github.com/iluuua}{github.com/iluuua} \textperiodcentered{}
  +7 995 410-33-10
\end{center}

% ===== Summary =====
\section{О себе}
Пишу бэкенд на Python. Бэкенд B2B SaaS-продукта ReachFlow спроектировал с нуля и довёл до продакшена. Он собран
на FastAPI поверх asyncpg и PostgreSQL, с фоновыми воркерами на Redis и Dramatiq и асинхронными
интеграциями с Telegram по MTProto и с языковыми моделями. Занимаюсь научной работой, написал статью
по молекулярно-динамическому моделированию алюминия с включением Fe$_4$Al$_{13}$, подробности на GitHub.
Алгоритмическая база олимпиадная, на C++.

% ===== Skills =====
\section{Навыки}
\begin{itemize}
  \item \textbf{Бэкенд.} Python, FastAPI, Pydantic, asyncio, REST API, JWT, OAuth
  \item \textbf{Данные.} PostgreSQL, SQL, asyncpg, pgvector, Redis
  \item \textbf{Очереди и воркеры.} Dramatiq, фоновые задачи, отложенные и запланированные действия
  \item \textbf{Инфраструктура.} Docker, Docker Compose, nginx, Kubernetes, CI/CD, Linux
  \item \textbf{Качество и наблюдаемость.} pytest, unittest, ruff, health-чеки, структурные логи, request-id
        и trace-id, OpenTelemetry
  \item \textbf{Интеграции.} Telegram API, MTProto, Telethon, Google OAuth, LLM через OpenRouter
  \item \textbf{Алгоритмы.} C++, алгоритмы и структуры данных
\end{itemize}

% ===== Experience =====
\section{Опыт работы}

\role{ReachFlow, Python backend и product engineer}{с января 2026}
\textit{B2B SaaS для автоматизации продаж в Telegram \textperiodcentered{} \href{https://reachflow.tech}{reachflow.tech}}
\begin{itemize}
  \item Спроектировал и написал многотенантный бэкенд на FastAPI и asyncpg. В нём 14 доменных роутеров, среди них
        Auth, Workspaces, CRM, Inbox, Discovery, Outreach, Sequences, AI~Brain, Billing, Connections, Import и Stats.
        Вход идёт по JWT и через Google OAuth с проверкой подписи по JWKS, а данные разных воркспейсов изолированы
        и на уровне запросов, и RLS-политиками Postgres.
  \item Спроектировал схему PostgreSQL. Сейчас в ней 8 доменных схем, 69 таблиц, 168 индексов и 33 forward-миграции,
        плюс bootstrap базы, диагностические эндпоинты и smoke-проверки схемы.
  \item Сделал фоновую обработку на Redis и Dramatiq. Отдельные акторы отвечают за диалог и за планировщик, входящие
        сообщения батчатся, follow-up и запланированные действия откладываются во времени, а очереди восстанавливаются
        после сбоя.
  \item Подключил Telegram по MTProto через Telethon. Написал login-flow, хранение и восстановление session string,
        работу через MTProxy и SOCKS, синхронизацию входящих и исходящих, свои лимиты отправки и обработку
        flood-ограничений.
  \item Построил слой оркестрации LLM поверх OpenRouter. В него входят профили моделей и цепочки фолбэков, классификация
        состояния лида, извлечение фактов, выжимка диалога для менеджера и трейсы вызовов, а эмбеддинги живут
        в pgvector.
  \item Довёл сервис до продакшена. Docker и Docker Compose, nginx обратным прокси с тем же origin для \texttt{/api},
        health-чеки трёх уровней вплоть до валидации схемы БД, структурные логи с request-id и trace-id
        по OpenTelemetry, CI на GitHub Actions с тестами и ruff, 86 тестовых модулей. Деплой идёт на Amvera
        с managed PostgreSQL, для воркеров написаны Kubernetes-манифесты.
\end{itemize}
\vspace{3pt}

\role{Фриланс, Python-разработчик}{с июля 2024}
\begin{itemize}
  \item Телеграм-бот для дистрибьютора пластиковых окон. Асинхронный бот на Aiogram, интеграция с Битрикс24
        по REST API, PostgreSQL с векторным поиском для авто-консультаций, уведомления о смене статуса заказа.
        По оценке заказчика бот автоматизировал около 80\% типовых консультаций.
  \item Телеграм-бот для школы (\texttt{@shkolo\_teacher\_bot}). Внутри три роли, ученик, учитель и админ,
        автоматизация расписания, новостей и отчётности, деплой в Docker.
\end{itemize}
\vspace{3pt}

\role{Near AI, разработчик решений на C++}{с июня 2023 по июль 2024}
\begin{itemize}
  \item Решал олимпиадные алгоритмические задачи на C++ для обучения децентрализованного ИИ. Оптимизировал
        по времени и памяти, оценивал асимптотику, рефакторил.
  \item Оформлял решения в LaTeX и ревьюил код других участников, команда была международная.
\end{itemize}

% ===== Education =====
\section{Образование}
\textbf{Московский политехнический университет}, прикладная математика и информатика, выпуск 2029.
\textbf{НИУ ВШЭ, ФКН}, вольнослушатель курсов направления <<Прикладная математика и информатика>>.
В 2025 году прошёл Т-Поколение <<Алгоритмы>> в параллели B, Deep Learning School и научную школу ОЦ <<Сириус>>.

% ===== Additional =====
\section{Дополнительно}
GitHub \href{https://github.com/iluuua}{github.com/iluuua} \textperiodcentered{} английский C1
\textperiodcentered{} русский родной

\end{document}
